\section{Direct $n=2$ matter-growth and Weyl transfer}
\label{sec:r1-n2-growth-lensing-transfer}

The reduced physical $n=2$ propagator can be coupled directly to minimally coupled matter through covariant energy--momentum conservation, without invoking the curvature-defective Newtonian-gauge DAE.

For pressureless matter in unitary gauge,
\begin{align}
\dot v_m&=-\delta n,\\
\dot\delta_m+3\frac{d}{dt}(Hv_m)&=-\left(3\dot\zeta+\frac{k_2^2}{a^2}\chi\right)+\frac{k_2^2}{a^2}v_m,
\end{align}
with $k_2^2/a^2=8K/a^2$ for the physical manuscript $n=2$ harmonic.

In the same unitary-gauge convention the gauge-invariant potentials are
\begin{equation}
\Psi=\delta n+\dot\chi,\qquad\Phi=\zeta+H\chi,
\end{equation}
so the local Weyl/lensing combination in this sign convention is $\Psi-\Phi$.

The response is integrated for both basis columns of the physical state $(\zeta,d\zeta/dN)^T$, with zero matter response at the $z=2.1$ anchor.  Thus the result is a row-vector kernel on the same two-dimensional physical state space.

\subsection{Present-epoch state-space kernels}
\begin{align}
\mathcal F^{\rm state}_{f\sigma_8}(0)&=(2.172303247,\,3.071170118\times10^{-2}),\\
\mathcal F^{\rm state}_{\rm Weyl}(0)&=(6.167895805\times10^{-2},\,-2.498933170\times10^{-2}).
\end{align}

These are not two new amplitudes.  They are the two components of linear functionals acting on the same initial $n=2$ state.  One amplitude gives unique predictions only after the normalized state direction is independently fixed; the distance row alone has rank one.

\subsection{Claim boundary}
The $f\sigma_8$ object above is a state-space fractional response kernel, using the R1 baseline $dD/dN$ normalization.  It is not yet the final single-number directional prediction because the full $D_L$ calibration functional has not yet selected the initial-state vector.
Likewise $\Psi-\Phi$ is the local Weyl potential combination; the integrated convergence/shear line-of-sight kernel remains to be evaluated.
